\section{Experimental Setup}
\label{sec:setup}

\subsection{Task Corpus \& Stratified Template Splits}
We evaluate LASR-MCP on a synthetic enterprise-style supply-chain tool execution benchmark (\texttt{create\_purchase\_order}) governing multi-item warehouse replenishment. The benchmark corpus comprises 80 pre-specified task instances partitioned strictly across syntactic template families:
\begin{itemize}
    \item \textbf{Training Tasks ($N=48$)}: Derived from template families \texttt{family\_a} and \texttt{family\_b}, used to train the frozen semantic predictor.
    \item \textbf{Evaluation Tasks ($N=32$)}: Drawn exclusively from \texttt{family\_c} with distinct lexical and syntactic structures (serving as unseen held-out test tasks for $H_1$ intent prediction, and development/validation tasks for $H_2$--$H_4$ systems evaluation).
\end{itemize}
The 32 evaluation tasks span four diverse operational archetypes ($N=8$ tasks each):
\begin{enumerate}
    \item \textbf{C1: Routine Regional Replenishment}: High-frequency, standard priority orders governed by basic warehouse routing.
    \item \textbf{C2: Expedited Emergency Orders}: Critical priority dispatches with accelerated freight logistics and rigid budget limits.
    \item \textbf{C3: Specialized Bulk Procurement}: Large-quantity orders where queries omit explicit warehouse locations, testing model behavior under unstated default enums.
    \item \textbf{C4: Multi-Item Depot Rebalancing}: Complex multi-line consignments enforcing nested \texttt{maxItems} and item-level SKU validation.
\end{enumerate}

\subsection{Controlled Experimental System Arms}
To compare the contributions of the schema-resolution components in Equation~\eqref{eq:three_tier}, we evaluate six controlled system arms under identical prompt conditions:
\begin{itemize}
    \item \textbf{Full Schema Control (\texttt{full})}: Static exposure of the monolithic 552-token schema upfront.
    \item \textbf{Invariant Baseline (\texttt{ci\_k0})}: Invariant Foundation alone ($C_I^* = 339$ tokens) without task-conditioned prediction or reactive recovery.
    \item \textbf{Proactive Static Control (\texttt{ci\_k1\_static})}: Invariant Foundation + Top-1 LASR predicted constraint without reactive recovery.
    \item \textbf{Proactive-Reactive LASR (\texttt{ci\_k1\_reactive})}: Invariant Foundation + Top-1 LASR predicted constraint with dynamic in-tree recovery.
    \item \textbf{Proactive-Reactive LASR (\texttt{ci\_k2\_reactive})}: Invariant Foundation + Top-2 LASR predicted constraints with dynamic in-tree recovery.
    \item \textbf{Model-Conditional Reference (\texttt{oracle\_reactive})}: Invariant Foundation + reference constraint set $O(x) = C_I^* \cup C_T^*(x)$ with reactive recovery. $O(x)$ was constructed from task specifications prior to execution and validated against the schema validator.
\end{itemize}

\subsection{Dual Evaluation Criteria \& Metric Definitions}
Submissions are evaluated on two orthogonal metrics alongside token efficiency:
\begin{enumerate}
    \item \textbf{Structural Schema Validity (\%)}: Conformance to JSON Schema Draft 2020-12 verified by an independent automated validator.
    \item \textbf{Semantic Task Correctness (\%)}: Exact domain logic validation verifying that extracted line items, quantities, priority codes, and supplier IDs satisfy user intent. Evaluation is strictly hierarchical: semantic correctness requires structural validity as a precondition ($\text{First-Pass Hierarchical} = \text{Structural} \land \text{Semantic}$), bounding semantic correctness by structural success.
    \item \textbf{Recorded Token Metric \& Accounting Reconciliation}: The primary token metric is defined strictly as the provider-recorded prompt-plus-candidate token subtotal per recorded attempt:
    \begin{equation}
    \text{Recorded } P+C = \text{prompt\_tokens} + \text{completion\_tokens}
    \end{equation}
    directly derived from Google's \texttt{usageMetadata} (\texttt{promptTokenCount} + \texttt{candidatesTokenCount}, which identically equals provider \texttt{totalTokenCount} across all 7,860 calls in our benchmark dataset as thinking and client tool prompt tokens are zero). In earlier preliminary benchmark drafts, heuristic offline schema token estimates ($S_0 + \Delta S$) were inadvertently added to provider prompt tokens, creating an artificial double-counting distortion that heavily penalized Full Control (whose monolithic schema is already serialized inside the prompt). In v2.0.3, we eliminate this double-counting artifact entirely: heuristic schema sizes ($S_0, \Delta S$) are decoupled and retained strictly as offline structural diagnostics, while all economic and token workload evaluations rely exclusively on provider-reported $P+C$ token usage from \texttt{usageMetadata}.
\end{enumerate}

\subsection{Execution Environment \& Live Telemetry Ledger}
All live experiments were executed against model ID \texttt{gemini-3.5-flash-lite} hosted on Google Cloud Generative Language API via direct REST HTTP endpoints. The provider telemetry explicitly records \texttt{requested\_model} and \texttt{provider\_resolved\_model\_version} as \texttt{gemini-3.5-flash-lite} (returned in the API \texttt{modelVersion} field); the API reported this model alias directly, and a separate immutable backend snapshot tag was not exposed by the endpoint. In our client configuration, \texttt{LiveLLMClient} operated at temperature $0.2$, and requests transmitted a fixed execution seed (\texttt{seed=42}). Across the 20 repeated executions per task-arm combination, the model produced deterministic identical first responses and outcomes under the repeated fixed seed; the 32 distinct task IDs represent the paired resampling units of inference, and the 20 replicates verify execution determinism under fixed decoding conditions rather than representing independent stochastic draws. 

The live empirical dataset comprises exactly $N=3,840$ completed trials across the 32 tasks and 6 experimental arms, producing $7,860$ response-bearing provider calls (all HTTP 200, with 100\% complete provider \texttt{usageMetadata} recording \texttt{promptTokenCount}, \texttt{candidatesTokenCount}, and \texttt{totalTokenCount}). Transient network rate limits (HTTP 429 and 503) triggered exactly 154 client transport retries; these transport-level backoffs are cleanly isolated and reported separately from logical agent attempts. 

We distinguish the full multi-session data collection timeline from the duration of the final resumed session: the total collection spanned multiple days across paid and quota-limited API sessions, completing 3,149 trials in initial sessions followed by a final resumed session completing the remaining 691 trials in approximately 2.8 hours. Operational execution ledgers are preserved in the reproducibility archive.

\subsection{Statistical Inference: Paired Task-Cluster Percentile Bootstrap}
For all comparative workload evaluations between experimental arms, we employ a paired task-cluster percentile bootstrap over the 32 task clusters. The 32 distinct task IDs constitute the resampling units of inference. For each bootstrap iteration ($B=10,000$ resamples, deterministic seed 42), task clusters are sampled with replacement, computing paired workload differences and cluster-level ratios. Confidence intervals are reported as empirical 2.5\% and 97.5\% percentiles of the bootstrap distribution.

\subsection{Evaluation Hygiene \texorpdfstring{\&}{and} Benchmark Scope}
To maintain strict methodological hygiene, we explicitly delineate the empirical role of each benchmark split:
\begin{itemize}
    \item \textbf{Held-Out Family Generalization for $H_1$}: The intent-conditioned predictor is implemented as a custom multi-label log-odds-weighted lexical classifier with class priors and smoothed IDF weights. It was trained on 48 tasks (\texttt{family\_a}, \texttt{family\_b}) and evaluated on 32 tasks (\texttt{family\_c}) drawn from a held-out syntactic template family. While template framing structures are strictly partitioned, shared enterprise procurement vocabulary exists across splits (103 of 149 test vocabulary words, $69.1\%$ overlap), and all four target constraint combinations recur across train and test families. Predictor weights were formally frozen prior to all systems evaluations.
    \item \textbf{Primary Frozen Benchmark for $H_2$--$H_4$}: The invariant foundation $C_I^* = 339$ tokens was calibrated on an 8-task calibration probe (\texttt{BLIND-C1-013}, \texttt{BLIND-C1-014}, \texttt{BLIND-C2-013}, \texttt{BLIND-C2-014}, \texttt{BLIND-C3-013}, \texttt{BLIND-C3-014}, \texttt{BLIND-C4-013}, \texttt{BLIND-C4-014}) sampled directly from the \texttt{family\_c} task distribution (two tasks per cluster). Consequently, the 8-task probe overlaps the 32 evaluation tasks; we explicitly disclose this tuning/evaluation overlap and evaluate $H_2$--$H_4$ as characterizing systems behavior within this primary frozen benchmark corpus rather than serving as an untouched, blind end-to-end generalization test.
\end{itemize}

\subsection{Data Provenance \texorpdfstring{\&}{and} Trace Replay Verification}
All benchmark code, traces, and frozen datasets are archived at \url{https://github.com/AnandSuchak/mcp-inventory-bench}. While live API re-executions may exhibit variance from hosted model updates, recorded execution traces support deterministic replay of the reported experiments, verified by automated cryptographic manifests:
\begin{center}
\footnotesize
\textbf{Cryptographic Artifact Provenance:}\\
\textit{Frozen Benchmark ($N=32$):}\\
\texttt{\scriptsize SHA-256: 78c2c88ed5128e0f296fa29e0c258589}\\
\texttt{\scriptsize \phantom{SHA-256: }b516f3adb14b0b8da491fa73bc3ffc11}\\
\textit{v2.0.3 Live Benchmark Candidate ($N=3,840$):}\\
\texttt{\scriptsize SHA-256: fc82d33b8992914ee130ef2805427ef2}\\
\texttt{\scriptsize \phantom{SHA-256: }a4f01541243e0f93be120e2d20ae7f5b}\\
\texttt{\scriptsize Commit: 9807340abc38170e0090a102547ee4309c01d78b}
\end{center}
Automated regression test suites ensure complete data immutability.

